\documentclass[UTF8,a4paper,10pt]{ctexart}
\usepackage[a4paper,margin=1.75cm]{geometry}
\usepackage{amsmath,amssymb,booktabs,tabularx}
\usepackage[most]{tcolorbox}
\setlength{\parindent}{2em}\setlength{\parskip}{0.3em}
\newtcolorbox{corebox}[1]{colback=gray!8,colframe=black!70,title=\textbf{#1},boxrule=.7pt,arc=1mm}
\begin{document}
\section*{B07\quad 随机变量变换与 Jacobian：概率质量守恒}
\textbf{依赖：}B02+B06A/B06B。只解释变换方向、支持迁移和质量守恒。
\begin{tcolorbox}[title=母接口]
原坐标区域的概率质量不变：
\[
f_{U,V}(u,v)=f_{X,Y}(x(u,v),y(u,v))
\left\lvert\frac{\partial(x,y)}{\partial(u,v)}\right\rvert.
\]
\end{tcolorbox}
\subsection*{1.\ 一一变换}
先写 $(u,v)=T(x,y)$ 的逆像 $(x,y)=T^{-1}(u,v)$ 与新支持域，再使用逆变换 Jacobian。矩阵和行列式方向必须与“密度乘微元保持不变”
\[
f_{X,Y}\,dxdy=f_{U,V}\,dudv
\]
一致。
\subsection*{2.\ 一维与多分支}
若 $Y=g(X)$ 在目标区间单调，
\[
f_Y(y)=f_X(g^{-1}(y))\left\lvert\frac{d}{dy}g^{-1}(y)\right\rvert.
\]
非单调时按所有满足 $g(x_i)=y$ 的分支求和，并分别检查原支持。
\subsection*{3.\ 失败边界}
\begin{itemize}
\item 只换公式不换 support 会产生错误概率；
\item Jacobian 为零或变换非光滑时不能机械套公式；
\item 多对一映射必须分支求和，不能挑一个逆函数。
\end{itemize}
\subsection*{4.\ 最小例子与调用}
若 $X\sim U(-1,1)$ 且 $Y=X^2$，对 $0<y<1$ 有两个逆像 $\pm\sqrt y$，故
\[
f_Y(y)=\frac12\frac{1}{2\sqrt y}+\frac12\frac{1}{2\sqrt y}
=\frac{1}{2\sqrt y}.
\]
题面出现非单调变换时，先列完整原像与新支持，再对所有分支求和。
\begin{corebox}{检查句}
守恒等式中的微元方向是什么？新域是否是完整逆像？是否存在多个分支、奇异点或边界缺口？
\end{corebox}
\end{document}
